\section*{Seguridad en el Manejo de Materiales Radiactivos}

\textbf{Instrucciones Generales de Seguridad:} \\
Es conveniente utilizar un monitor de radiación durante la operación, para tener la seguridad de que la fuente está blindada de forma apropiada y para verificar que los límites de exposición no se excedan.

Por ningún motivo se permite el traslado del equipo ensamblado para hacer nuevas exposiciones en tramos largos o cambiar de nivel de posición con el equipo ensamblado; se deberá desensamblar y mantener sus equipos de monitoreo encendidos hasta que la fuente guía esté completamente blindada dentro del contenedor. Al término de cada jornada de trabajo el personal ocupacionalmente expuesto deberá asegurarse de que la fuente quede bien blindada dentro del contenedor, revisando los niveles de radiación que presenta y comparándolos con los que tenía antes de iniciar la jornada.

Al término de la jornada el material deberá ser guardado en el Almacén de Material Radiactivo de nuestra compañía. Cuando por necesidades de servicio el contenedor radiactivo se encuentre fuera de la ciudad donde se localiza nuestro almacén de Material Radiactivo, deberá mantenerse en el portacontenedor de la unidad y guardarse en un estacionamiento cerrado y con vigilancia, por ningún motivo se deberá dejar la unidad estacionada en la vía pública con el material radiactivo en ella y sin vigilancia.

% Sección sobre almacenamiento y transporte
\textbf{Almacenamiento y Transporte:} \\
Respecto al almacenamiento y transporte del material radiactivo dentro del territorio nacional, este es regido por las disposiciones y normas establecidas por la CNSNS (\textit{Comisión Nacional de Seguridad Nuclear y Salvaguardias}).

La unidad en que se transporte el material radiactivo se le debe colocar un rótulo o etiqueta en cada uno de sus lados. En la mitad superior ira el símbolo internacional que indica la presencia de radiación y el color negro. En la mitad inferior del rótulo se indica la leyenda ``NO PERMANEZCA INNECESARIAMENTE EN CERCA DE ESTE VEHÍCULO'' con colores que hagan contraste.

Además, el vehículo deberá estar provisto de un cordón amarillo, necesario para acordonar dos veces el perímetro de dicha unidad, postes para colocar el cordón, así, como letreros con el símbolo internacional de que indica la presencia de radiación.

El tránsito de un vehículo que contenga material radiactivo se hará por la ruta más corta y siempre observando los límites de velocidad señalada por las autoridades.